% =====================================================================
%  Diploma in Data Science and Artificial Intelligence Applied to
%  Health Research - Mentorship Project 3
%  EEG and motor imagery for the rehabilitation of stroke patients
%
%  English version of the submitted report (../../Informe.pdf),
%  translated section by section.
%
%  Compile with pdfLaTeX (in Overleaf: Recompile).
% =====================================================================
\documentclass[12pt,a4paper]{article}

\usepackage[utf8]{inputenc}
\usepackage[T1]{fontenc}
\usepackage{mathptmx}
\usepackage[margin=2.54cm]{geometry}
\usepackage{setspace}
\usepackage{enumitem}
\usepackage{microtype}
\usepackage[hidelinks]{hyperref}
\hypersetup{
  pdftitle={EEG and Motor Imagery for the Rehabilitation of Stroke Patients},
  pdfauthor={Mariel Anachuri, Clara Inés Chungara, Camila Giolitto}
}

\setstretch{1.15}
\setlength{\parindent}{1.25cm}
\setlength{\parskip}{0pt}
\pagestyle{empty}

% Section headings in the style of the original report
\newcommand{\mainsection}[1]{%
  \par\bigskip\noindent\textbf{#1}\par\nopagebreak}
\newcommand{\subheading}[1]{%
  \par\noindent\underline{#1}\par\nopagebreak}

\setlist[itemize]{leftmargin=1.9cm, itemsep=0pt, topsep=0pt, parsep=0pt}
\setlist[enumerate]{leftmargin=1.9cm, itemsep=0pt, topsep=0pt, parsep=0pt}

\begin{document}

% ---------------------------------------------------------------------
%  Title block
% ---------------------------------------------------------------------
\begin{center}
  {\textbf{EEG and Motor Imagery for the Rehabilitation of Stroke Patients}}\\[2pt]
  {\textbf{Feature Engineering and Per-Subject Training on the Cho2017 Dataset}}\\[2pt]
  Mariel Anachuri, Clara Inés Chungara, Camila Giolitto\\[10pt]
  {\small\textbf{\textit{Diploma in Data Science and Artificial Intelligence Applied to Health Research}}}\\[2pt]
  \textit{Mentorship Project 3 --- Mentor: Bruno Zorzet}
\end{center}
\vspace{-4pt}
\noindent\rule{\textwidth}{0.4pt}

% ---------------------------------------------------------------------
%  Abstract
% ---------------------------------------------------------------------
\noindent\textbf{Abstract}\par
In this work, an EEG signal processing and machine learning pipeline was developed to classify
two motor imagery tasks (left hand and right hand) using the Cho2017 database. A feature
engineering stage was carried out on the EEG signal, transforming the original recordings into
a tabular dataset by means of frequency, temporal and spatial descriptors. Machine learning
models were trained on this feature set under a \textit{within-subject} strategy, comparing
Decision Tree, Random Forest, SVM and CSP+LDA as the baseline model. The evaluation made it
possible to analyse the performance of the classifiers and the influence of the individual
variability associated with the \textit{BCI illiteracy} phenomenon.

\noindent\textbf{Keywords}\par
EEG, MI-BCI, Feature Engineering, Machine Learning, BCI Illiteracy

% ---------------------------------------------------------------------
\mainsection{I. INTRODUCTION AND PROBLEM STATEMENT}
% ---------------------------------------------------------------------
Motor imagery brain-computer interfaces (\textit{Motor Imagery Brain-Computer Interfaces},
MI-BCI) make it possible to decode a user's movement intention from electroencephalography
(EEG) signals, which makes them a tool of interest for motor rehabilitation applications in
patients who have had a cerebrovascular accident (\textit{stroke}).

The deployment of these technologies faces two highly complex challenges:
\begin{enumerate}
  \item High inter-subject variability: EEG signals are inherently non-stationary and show
        spectral and spatial signatures that are unique to each individual, which prevents the
        efficient use of global classification models.
  \item The BCI illiteracy phenomenon (or BCI inefficiency): it is statistically estimated that a
        significant percentage of users ($\sim$20\% in the general literature) are unable to
        produce brain responses that are sufficiently discriminative or voluntary to operate a
        BCI system reliably.
\end{enumerate}

To address this problem, this work uses the Cho2017 database (52 subjects, $\sim$200 trials per
subject, balanced across two classes), reducing the analysis space to only 3 critical channels:
C3, Cz and C4, resampled to 128 Hz.

The central objective of this project is to design, implement and critically evaluate a signal
processing and multi-domain feature engineering pipeline, comparing classical tabular algorithms
against the industry's spatial standard (CSP+LDA) under a realistic within-subject validation
scheme that quantifies the impact of BCI illiteracy.

% ---------------------------------------------------------------------
\mainsection{II. METHODOLOGY}
% ---------------------------------------------------------------------
A within-subject approach was adopted to mitigate the inter-individual variability of the EEG.
Starting from the preprocessed signals, represented as matrices of dimensions
trials $\times$ 3 channels (C3, Cz, C4) $\times$ time, the following stages were carried out:

\subheading{1. Extraction and construction of the feature set}
The target variable (motor intention) was coded in binary form (0: left hand, 1: right hand).
For each trial, 51 descriptors were generated, grouped into three families:
\begin{itemize}[label=-, leftmargin=3.4cm]
  \item Frequency: total power (RMS) and spectral power in 6 frequency bands (4--8, 8--12,
        12--16, 16--20, 20--24 and 28--30 Hz), using Welch's method (Welch, 1967).
  \item Temporal and complexity: kurtosis, skewness, Katz fractal dimension (Katz, 1988),
        autoregressive coefficients (order 2) and permutation entropy (order 3).
  \item Spatial and connectivity: cross-correlation, spectral coherence (MSC) and Phase Locking
        Value (PLV), between C3--Cz, C4--Cz and C3--C4.
\end{itemize}

Consolidating the 52 subjects resulted in a dataset of 10,520 trials $\times$ 54 columns
(51 features + 3 metadata columns: subject, coded label and its readable version).

\subheading{2. Exploratory analysis and data leakage prevention}
Boxplot + stripplot charts were generated per channel and per class to assess the overlap
between distributions. It was determined that univariate visual inspection is not enough to
anticipate the performance of a multivariate classifier; a valid conclusion requires formal
hypothesis tests (Shapiro-Wilk, followed by a parametric or non-parametric test). No feature
selection or transformation was applied before the split, to avoid data leakage.

\subheading{3. Machine learning pipeline and evaluation}
For each subject, the data were divided with a stratified 80/20 split, keeping the test set
only for the final evaluation. Within the 80\% training set, a GridSearchCV object coupled to a
stratified 5-fold cross-validation was applied. Feature scaling (StandardScaler), where
applicable, was strictly encapsulated inside a Scikit-Learn Pipeline. This ensured that the
normalisation parameters were estimated and applied using only the training data of each fold,
preventing validation or test information from biasing the fit of the models.

To close this workflow, the pipeline trained and optimised the four selected models
independently for each of the 52 subjects. At the end of this stage, a performance matrix was
consolidated containing, for each individual, four final accuracy metrics computed exclusively
on their own held-out test sets.

\subheading{4. Model Choice and Tuning Criteria}
\begin{itemize}[label=-, leftmargin=1.3cm]
  \item Model justification: the aim was to cover the full methodological spectrum, in order to
        contrast the paradigm of handcrafted features against a method that learns its own spatial
        projection of the data. Under this premise, the following were chosen: Decision Tree
        (interpretable, low-complexity base model), Random Forest (ensemble robust to noise;
        Breiman, 2001), SVM (maximum-margin model with non-linear capacity via kernel; Cortes \&
        Vapnik, 1995) and CSP + LDA (industry standard in MI-BCI; Ramoser et al., 2000), the
        latter fitted directly on the raw EEG signal of the 3 spatial channels and used as the
        mandatory baseline.
  \item Hyperparameter tuning: the configurations were optimised through the GridSearchCV grid,
        evaluating critical parameters such as depth and pruning in trees (max\_depth,
        min\_samples\_leaf), the number of estimators in Random Forest, the regularisation (C)
        together with the kernel type in SVM, and the number of spatial components in CSP. The
        final selection was based strictly on the mean accuracy obtained in cross-validation.
  \item Class balancing: since the experimental design of the Cho2017 database provides naturally
        balanced classes ($\sim$50/50 between both hands), it was decided to omit any resampling
        technique.
\end{itemize}

\subheading{5. Configuration of the BCI Illiteracy Filter}
To rigorously distinguish real predictive performance from statistical chance, given the finite
size of the test sample, an exclusion filter was designed based on modelling the Binomial
Distribution under the null hypothesis (H$_0$) of a purely random classifier.

The theoretical decisions behind the filter considered a theoretical probability of success of
$p = 0.5$ (due to the dichotomous class balance) and a significance level of $\alpha = 0.05$, to
tolerate a maximum error margin of 5\%. Considering the size of the test set, the binomial
distribution set a mathematical minimum of 0.650 accuracy to reject the chance hypothesis and
consider a subject ``eligible''. This official cut-off criterion was applied by strictly
evaluating the performance obtained with the CSP+LDA baseline model.

% ---------------------------------------------------------------------
\mainsection{III. ANALYSIS OF RESULTS AND DISCUSSION}
% ---------------------------------------------------------------------
\subheading{1. Global Comparison of Models (52 Subjects)}
All four models exceeded, on average, the theoretical chance level of 0.5 (where 0.5 represents
the probability of a purely random correct answer in binary classification). Support Vector
Machines (SVM) and Random Forest led the overall predictive performance, with mean accuracy
values of 0.677 and 0.665, followed by the Decision Tree (0.619) and CSP+LDA (0.591), which had
the lowest performance in the study. The high dispersion observed in the standard deviations
reflected the marked inter-subject variability, which supports the methodological decision to
train the models under the within-subject approach.

The low overall performance obtained by CSP + LDA (0.591) contrasts with the performance usually
reported in the classical literature for this dataset. The explanation for this may lie in the
hardware restriction implemented in this experimental design. The CSP algorithm is a spatial
optimiser that requires high dimensionality (many electrodes, such as the 64 of the original
Cho2017 montage) to compute robust cross-covariance matrices. By reducing the analysis space to
only 3 channels (C3, Cz, C4), CSP loses practically all of its spatial degrees of freedom and
becomes highly sensitive to noise.

In contrast, the tabular models showed that when the hardware limits spatial resolution,
multi-domain feature engineering manages to compensate for the loss of physical sensors, finding
complex non-linear decision boundaries (such as those of SVM with an RBF kernel) that a purely
spatial extractor does not detect on bare raw signals.

\medskip
\subheading{2. BCI Illiteracy Filter and Analysis of the Eligible Subpopulation}
Taking the reference classifier (CSP+LDA) as the baseline, only 14 of the 52 subjects (26.9\%)
qualified as eligible by exceeding the binomial cut-off threshold of 0.650, while 38 subjects
(73.1\%) remained below the threshold, under the condition of BCI illiteracy.

The statistical filter did not modify the data or require retraining the classifiers. Instead,
the rows corresponding to the 14 subjects declared eligible were isolated within the
consolidated results table, in order to evaluate the behaviour of the set of models under
validated neurophysiological profiles. When the metrics were computed exclusively on this
filtered subgroup, a drastic and systematic improvement of between 12 and 16 percentage points
was observed in the mean accuracy of all architectures:
\begin{itemize}[label=\textbullet, leftmargin=1.3cm]
  \item Decision Tree: rose from 0.619 to 0.739 (+0.120).
  \item CSP + LDA (Baseline): rose from 0.591 to 0.741 (+0.150).
  \item Support Vector Machine (SVM): rose from 0.677 to 0.824 (+0.147).
  \item Random Forest: rose from 0.665 to 0.821 (+0.156).
\end{itemize}

This quantitative jump statistically supports that the filter based on the binomial
distribution effectively separates two populations with qualitatively different bioelectrical
signatures. By removing the noise introduced by subjects unable to modulate their cortical
rhythms, the performance of all algorithms increased notably, including CSP+LDA itself
(+15 points), confirming that when the physiological signature exists, the spatial algorithm
works even with hardware limitations.

From a clinical perspective, SVM (0.824) and Random Forest (0.821) broke the critical barrier of
0.80 mean accuracy. In biomedical engineering and in practice, an accuracy above 0.80 is
established as the minimum operational threshold to consider a brain-computer interface
functional and usable by a patient in post-stroke rehabilitation, allowing them to control
neurofeedback software fluidly and without experiencing frustration from false commands.

\medskip
\subheading{3. Strengths, Limitations and Approaches Not Implemented}
\begin{itemize}[label=\textbullet, leftmargin=1.3cm]
  \item Strengths: the complete protection against data leakage, through the use of encapsulated
        Scikit-Learn pipelines, guaranteed a statistically honest evaluation. In addition, the
        superiority of SVM and Random Forest over CSP+LDA showed that when the hardware limits
        spatial resolution (only 3 channels), multi-domain feature engineering (Welch power, Katz
        complexity, PLV/MSC connectivity) is the most robust strategy to recover discriminative
        patterns of motor intention.
  \item Limitations: the sample size per subject ($\sim$200 trials) is limited for training
        traditional models with dense grid searches, which introduces latent risks of
        overfitting at the lower end, among difficult users. In addition, the PSD sub-bands were
        fixed globally and uniformly, ignoring that the central frequency of the mu and beta
        rhythms varies physiologically between individuals (and even more so after a stroke).
  \item Approaches and Ideas Not Implemented: due to computational time constraints and agreed
        changes of scope, the implementation of formal feature selection algorithms (SelectKBest
        via ANOVA, RFE or Lasso L1 regularisation) to compress and optimise the 51 features was
        left out of the final deliverable. Likewise, it was not possible to carry out the
        non-linear graphical projection of the high-dimensional space with algorithms such as
        t-SNE and UMAP, originally planned to visually map the three-dimensional geometric
        separation between the eligible population and the population with BCI illiteracy.
\end{itemize}

% ---------------------------------------------------------------------
\mainsection{IV. CONCLUSIONS AND FUTURE WORK}
% ---------------------------------------------------------------------
\medskip
\subheading{1. General Conclusions}
This work showed that the viability of a Brain-Computer Interface (MI-BCI) based on traditional
machine learning depends on the individual neurophysiology of the user. Through experimentation
with multiple algorithms (Decision Trees, Random Forest, SVM and the CSP+LDA spatial baseline),
it was confirmed that the models are able to extract discriminative patterns of motor intention,
but their overall effectiveness is conditioned by inter-subject variability.

The integration of a strict statistical threshold transformed the analytical approach of the
project: it showed that a predictive model must assess its ability to exceed a significance
limit against chance. Although the findings revealed a high prevalence of ``BCI illiteracy'' in
the general population under reduced hardware configurations, when the eligible subjects are
statistically isolated, the classifiers substantially increase their predictive performance,
reaching accuracies suitable for clinical settings.

\medskip
\subheading{2. Future Work and Advanced Approaches}
To mitigate the limitations found and reduce the BCI illiteracy rate, the natural next step is
to move towards deep learning architectures. The implementation of specialised convolutional
neural networks such as EEGNet is proposed. These architectures make it possible to apply
simultaneous spatial and temporal convolutions to discover complex representations directly from
the native signal, with the potential to recover hidden patterns that escape traditional
statistical descriptors.

\end{document}
